\documentclass[12pt,a4paper,oneside,openany]{jsbook}
\usepackage[bthesis]{slab} % 卒論
%\usepackage[mthesis]{slab} % 修論
%\usepackage[dthesis]{slab} % 博論

\usepackage{graphicx}
\usepackage{mediabb}
\usepackage{ascmac}
\usepackage{fancybox}
%\usepackage{bigfoot}

% プログラムのコード掲載用
\usepackage[final]{listings}
\lstset{                 %listingsの設定
numbers=left,            %行番号を左
numberstyle=\scriptsize, %
stepnumber=1,            %1行おきに行番号を
numbersep=1zw,           %ソースと行番号の間隔
lineskip=-0.5zw,         %行間隔 要調整
basicstyle=\scriptsize\ttfamily,     %ttfamily
fontadjust=false,
xleftmargin=20pt,
frame=tRBl,
framesep=5pt,
framexleftmargin=17pt,
}

\newcommand{\figref}[1]{図 \ref{#1}}
\newcommand{\tabref}[1]{表 \ref{#1}}
\newcommand{\secref}[1]{章 \ref{#1}}
\newcommand{\equref}[1]{式 \eqref{#1}}
\newcommand{\lstref}[1]{Listing \ref{#1}}

\title{卒論のタイトル}
\author{信大 太郎}
\authorintitle{信~~大~~~~~太~~郎} % 表示用
\date{令和 X 年 X 月}
\nendo{令和 X 年度}
\studentid{XXTXXXXX}
\advisor{佐藤 真平 准教授}
\affiliation{信州大学工学部 電子情報システム工学科}

\begin{document}

\frontmatter
\maketitle

\chapter{概要}
\thispagestyle{plain}
\label{abstract}

概要aaaaaがいようごあいうｆｄｊさｋｌｆｊｄｋさｌｆｊｌｋ；ああああああああ
ああああああああああああああああああああああああああああああああああああああ
ああああああああああああああああああああ。
ｆｄさｆｄさｆｄさｆｄｊｌｋさｆｊだ。

にだんらくめああああいいいいいいいいいいいいいいいいいいいいいいいいいいいいいいいいいい
９いｆぢ；さｊｆｋｄｌさｊｆｋｄｌさｊｆｄさ


$Test^{hoge}$

$Test_{hoge}$


\tableofcontents
\listoffigures
\listoftables

\mainmatter
% 章番号が決まったらファイル名に連番を振っても良いし，
% include せずにすべてをこのファイルに含めても良い．
\chapter{序論}
\thispagestyle{plain}
\label{introduction}

\section{研究の背景と目的}

背景と目的。

\section{本論文の構成}

本論文の構成。

\chapter{背景}
\thispagestyle{plain}
\label{background}

背景。

\section{背景}

背景。

\chapter{前提知識}
\thispagestyle{plain}
\label{preliminary}

前提知識
\chapter{関連研究}
\thispagestyle{plain}
\label{related-works}

関連研究
\chapter{提案手法}
\thispagestyle{plain}
\label{proposal}

提案手法
\chapter{評価}
\thispagestyle{plain}
\label{evaluation}

評価

\chapter{結論}
\thispagestyle{plain}
\label{conclusion}

\section{結論}
結論

\section{今後の課題}
今後の課題

%\appendix
%\chapter{付録A}
%\chapter{付録B}

\backmatter
\chapter{謝辞}
\thispagestyle{plain}

謝辞

\bibliographystyle{myjunsrt}
\bibliography{library}

\end{document}
